% MCLF Book
% Models of Curves and Valuations
%
% Copyright (c) 2026 Stefan Wewers and contributors.
%
% License: CC BY 4.0. See the file LICENSE in the repository root.
%
% This file is part of the open book project hosted at:
% https://github.com/MCLF/mclf-book
%
% The text is work in progress. For citation, please refer to a
% specific Git commit or release.



% Basic input and font setup.
% With a modern LaTeX installation, utf8 is usually the default,
% but keeping this explicit is harmless for pdfLaTeX.
\usepackage[utf8]{inputenc}
\usepackage[T1]{fontenc}

% Standard mathematical packages.
\usepackage{amsmath}
\usepackage{amssymb}
\usepackage{amsthm}
\usepackage{mathtools}
\usepackage{stix}
%\usepackage{unicode-math}

% Configurable labels for standard list environments.
\usepackage{enumitem}

% Commutative diagrams.
\usepackage{tikz-cd}
\newcommand{\cdpunkt}{\mathrlap{\, .}}
\tikzset{
    symbol/.style={
        draw=none,
        every to/.append style={
            edge node={node [sloped, allow upside down, auto=false]{$#1$}}}
    }
}
\usepackage{tikzit}
\input{basic.tikzstyles}

\usepackage{comment}

% Graphics.
\usepackage{graphicx}
\usepackage{float}
\usepackage{tcolorbox}

% Visible notes for mathematical issues that require human resolution.
\newtcolorbox{authorwarning}{
    colback=red!5!white,
    colframe=red!75!black,
    fonttitle=\bfseries,
    title=Author input required
}

% References and hyperlinks.
% Load hyperref near the end.
\usepackage[
    colorlinks=true,
    linkcolor=blue,
    citecolor=blue,
    urlcolor=blue
]{hyperref}

% Theorem environments.
\theoremstyle{plain}
\newtheorem{theorem}{Theorem}[chapter]
\newtheorem{proposition}[theorem]{Proposition}
\newtheorem{lemma}[theorem]{Lemma}
\newtheorem{corollary}[theorem]{Corollary}

\theoremstyle{definition}
\newtheorem{definition}[theorem]{Definition}
\newtheorem{example}[theorem]{Example}
\newtheorem{construction}[theorem]{Construction}
\newtheorem{assumption}[theorem]{Assumption}

\theoremstyle{remark}
\newtheorem{remark}[theorem]{Remark}
\newtheorem{remarks}[theorem]{Remarks}

% Number equations independently by chapter.
\numberwithin{equation}{chapter}

% Slightly more flexible spacing in tables.
\renewcommand{\arraystretch}{1.1}
